\documentclass{article}
\usepackage{csquotes}
\usepackage[a4paper]{geometry}
\usepackage{fancyhdr}
\pagestyle{fancy}
\lhead{The Yellow Wallpaper}
\rhead{April 2026}
\begin{document}
\section{The Yellow Wallpaper}
\emph{Charlotte Perkins Gilam}'s short story \emph{The Yellow Wallpaper}, released in 1892, is about women's mental health and independende during the late 19th century, especially in the context of the \emph{rest cure}.
 
It follows an unnamed narrator, her husband, John, a physician, Jennie, the housekeeper, and Mary, who looks after the narrators son.
 
\subsection{Summary}
The narrator and John move to an old house. John put the narrator under the rest cure, forbidding her from doing most things. She journals anyways. Her mental state gets worse, she starts seeing women behind her rooms ugly wallpaper. She tries to save them, going insane herself.
 
\subsection{The Rest Cure}
The \emph{rest cure} was a \textquote{therapy}, invented by a male doctor, applied amost only to women.
 
Patients were forced to stay in bed the whole day for weeks or months at a time, not allowed to see friends or family, not allowed to do any real mental or physical activity, often treated just like children.
 
This \textquote{therapy} was often related to a diagnosis of \textquote{female hysteria}; where usually male doctors diagnosed women as \textquote{hysteric} because of symptoms of depression or because they did not conform gender roles.
 
Charlotte Perkins Gilman underwent the rest cure herself, nearly ending in a breakdown. \emph{The Yellow Wallpaper} was written to convince doctors that the rest cure was cruel and wrong. 
 
\subsection{Plot}
The narrator and her husband rent a colonial mansion for three months while their own house is being renovated. The narrator is supposed to recover from \emph{temporary nervous depression}, as her husband diagnosed it, there.
 
John is away for most of the day, sometimes the night too. 
 
The narrator would like to work, to write, but is forbidden from doing so; she, the men agree, needs rest. She secretly writes her diary anyways.
 
She doesn't like her room, which was probably used as a nursery. With time she becomes more and more irritated with its shades of yellow and its patterns.
 
Relatives arrive for the Fourth of July, but the narrator is too exhausted and cries for most of it.
 
She would like to visit her cousins, but John forbids her from doing so.
 
She starts seeing a woman behind the wallpaper, as if trapped by the pattern itself. She wants to leave the room, but John doesn't take her seriously. Later she also becomes obsessed with the wallpapers color and smell.
 
She thinks that the woman behind the wallpaper breaks free at daytime; she sees said woman outside the house, beside the road, through her window.
 
She locks the door and creeps around the room herself.
 
She, together with the woman, she thinks, peels off the wallpaper. She locks herself in and ties herself up with a rope to make sure she's not moved out.
 
She thinks that she'll have to return behind the wallpaper at night herself, wants to creep around the floor during the day.
 
John comes it, sees his wife on the floor, and faints. She creeps over him. 
 
\end{document}